% edit: 20180417
%DIF LATEXDIFF DIFFERENCE FILE
%DIF DEL main_unformatted_old.tex   Thu Sep 13 18:26:06 2018
%DIF ADD main_unformatted_new.tex   Wed Sep 19 14:54:22 2018
%DIF 2a2
% edit: 20180831 %DIF > 
%DIF -------
\documentclass{article}

\usepackage{xr}
\externaldocument[S]{supporting_information}
\usepackage{amssymb}
\usepackage[utf8]{inputenc}
\usepackage{graphicx, subcaption}
\usepackage{caption}
\usepackage{booktabs}
%DIF 11c12
%DIF < \usepackage{cite}
%DIF -------
\usepackage[superscript]{cite} %DIF > 
%DIF -------

\usepackage{setspace}
\doublespacing

%DIF 16a17
\ProvideTextCommand{\DJ}{OT1}{\raisebox{0.25ex}{-}\kern-0.6em D} %DIF > 
%DIF -------
\renewcommand{\thefootnote}{\fnsymbol{footnote}}

\interfootnotelinepenalty=10000
\date{}
%DIF PREAMBLE EXTENSION ADDED BY LATEXDIFF
%DIF UNDERLINE PREAMBLE %DIF PREAMBLE
\RequirePackage[normalem]{ulem} %DIF PREAMBLE
\RequirePackage{color}\definecolor{RED}{rgb}{1,0,0}\definecolor{BLUE}{rgb}{0,0,1} %DIF PREAMBLE
\providecommand{\DIFadd}[1]{{\protect\color{blue}\uwave{#1}}} %DIF PREAMBLE
\providecommand{\DIFdel}[1]{{\protect\color{red}\sout{#1}}}                      %DIF PREAMBLE
%DIF SAFE PREAMBLE %DIF PREAMBLE
\providecommand{\DIFaddbegin}{} %DIF PREAMBLE
\providecommand{\DIFaddend}{} %DIF PREAMBLE
\providecommand{\DIFdelbegin}{} %DIF PREAMBLE
\providecommand{\DIFdelend}{} %DIF PREAMBLE
%DIF FLOATSAFE PREAMBLE %DIF PREAMBLE
\providecommand{\DIFaddFL}[1]{\DIFadd{#1}} %DIF PREAMBLE
\providecommand{\DIFdelFL}[1]{\DIFdel{#1}} %DIF PREAMBLE
\providecommand{\DIFaddbeginFL}{} %DIF PREAMBLE
\providecommand{\DIFaddendFL}{} %DIF PREAMBLE
\providecommand{\DIFdelbeginFL}{} %DIF PREAMBLE
\providecommand{\DIFdelendFL}{} %DIF PREAMBLE
%DIF END PREAMBLE EXTENSION ADDED BY LATEXDIFF

\begin{document}

\title{Rapid and quantitative de-\emph{tert}-butylation for poly(acrylic acid) block copolymers and
influence on relaxation of thermoassociated transient networks}%$^\dag$}
\maketitle
\noindent\large{Alexei D. Filippov\textit{$^a$}, Ilse A. van Hees$^a$, Remco Fokkink$^a$, Ilja Voets\textit{$^b$} and Marleen Kamperman\textit{$^{a\ast}$}} \\%Author names go here instead of "Full name", etc.

\footnotetext{\textit{$^{a}$~Laboratory of Physical Chemistry and Soft Matter, Wageningen University \& 
Research, Stippeneng 4, 6708WE Wageningen, The Netherlands }}
\footnotetext{\textit{$^{b}$~Institute for Complex Molecular Systems, Eindhoven University of Technology,
P.O. box 513, 5600 MB Eindhoven, The Netherlands}}
\footnotetext{\textit{$^\ast$~Corresponding author, Tel: +31 317 482 358; E-mail: marleen.kamperman@wur.nl}}

\section*{Abstract}
The synthesis of charged polymers often requires the polymerization of protected monomers, followed by 
a polymer-analogous reaction to the polyelectrolyte product. We present a mild, facile method 
to cleave \emph{tert}-butyl groups from poly(\emph{tert}-butyl acrylate) blocks that yields
poly(acrylic acid) (pAA) blocks free of traces of the ester. The reaction utilizes a slight excess of HCl
in hexafluoroisopropanol (HFIP) at room temperature, and runs to completion within four hours. We compare deprotection
in HFIP with the common TFA/DCM method, and show that the latter does not yield clean pAA. We show the effect
of complete \emph{tert}-butyl cleavage on a ABA triblock copolymer, where poly(\emph{N}-
isopropylacrylamide) (pNIPAM) is A and pAA is B, by means of viscosimetry, DLS and SAXS 
on solutions above overlap. The pNIPAM blocks dehydrate, and their increased self-affinity above the 
lower critical solution temperature (LCST) results in network formation by the triblocks.
This manifests itself as an increase in viscosity and a slowing-down of the first-order correlation
function in light scattering. However, this stickering effect manifests itself exclusively when 
the pAA block is \emph{tert}-butyl-free. Additionally, SAXS shows that the conformational 
properties of \emph{tert}-butyl free pAA copolymers are markedly different from those with 
residual esters. Thus, we illustrate a surprising effect of hydrophobic impurities that acts
across blocks, and assert the usefulness of HCl/HFIP in pAA synthesis.\\

\section*{Introduction}
Acrylic acid (AA) is the smallest unsaturated carboxylic acid, and can be polymerized radically. By virtue of the 
carboxylic moiety, the resulting polymers are polyelectrolytes and have a pH-dependent charge and H-bonding capacity. 
In copolymers, AA is used to tailor responsive behavior of other polymeric building blocks, 
such as the swelling of hydrogels \cite{Lue2011, Yoo2000} and the lower critical solution temperature (LCST) 
of poly(\emph{N}-isopropylacrylamide) (pNIPAm) and its copolymers. \cite{Bokias2000} pAA is one of the most studied
polyanions in polyelectrolyte complexes, owing to its ability to form complexes with many polycations. 
\cite{Blocher2017, Liu2017}

In radical polymerizations, acrylic acid is unusually inclined to transfer reactions to solvent\DIFdelbegin \DIFdel{\mbox{%DIFAUXCMD
\cite{Loiseau2003a}
}%DIFAUXCMD
,}\DIFdelend \DIFaddbegin \DIFadd{,\mbox{%DIFAUXCMD
\cite{Loiseau2003a}
}%DIFAUXCMD
}\DIFaddend transfer to polymer\DIFdelbegin \DIFdel{\mbox{%DIFAUXCMD
\cite{Lena2017}
}%DIFAUXCMD
,}\DIFdelend \DIFaddbegin \DIFadd{,\mbox{%DIFAUXCMD
\cite{Lena2017}
}%DIFAUXCMD
}\DIFaddend and formation of Michael dimers\DIFdelbegin \DIFdel{\mbox{%DIFAUXCMD
\cite{LLauro2004}
}%DIFAUXCMD
.}\DIFdelend \DIFaddbegin \DIFadd{.\mbox{%DIFAUXCMD
\cite{LLauro2004}
}%DIFAUXCMD
}\DIFaddend These side reactions result in diminished control 
over linearity, molecular weight distribution, and end-group fidelity of the polymer products\DIFdelbegin \DIFdel{\mbox{%DIFAUXCMD
\cite{Maniego2013, LLauro2004}
}%DIFAUXCMD
.}\DIFdelend \DIFaddbegin \DIFadd{.\mbox{%DIFAUXCMD
\cite{Maniego2013, LLauro2004}
}%DIFAUXCMD
}\DIFaddend There are many
reports that detail pAA synthesis with the help of controlled radical polymerization (CRP) techniques, in particular reversible addition-fragmentation 
chain transfer \DIFdelbegin \DIFdel{\mbox{%DIFAUXCMD
\cite{Pelet2012, Loiseau2003a, Chaduc2013}
}%DIFAUXCMD
}\DIFdelend (RAFT).\DIFaddbegin \DIFadd{\mbox{%DIFAUXCMD
\cite{Pelet2012, Loiseau2003a, Chaduc2013}
}%DIFAUXCMD
}\DIFaddend Chain transfer agents are known to reduce branching and termination by ensuring 
a low concentration of active radicals\DIFdelbegin \DIFdel{\mbox{%DIFAUXCMD
\cite{Lena2017}
}%DIFAUXCMD
.}\DIFdelend \DIFaddbegin \DIFadd{.\mbox{%DIFAUXCMD
\cite{Lena2017}
}%DIFAUXCMD
}\DIFaddend Nonetheless, it is difficult to prepare strictly linear polymers directly from AA
due to its high rate and extensive repertoire of side reactions, and branching is often overlooked in reports on \DIFdelbegin \DIFdel{poly(acrylic acid)}\DIFdelend \DIFaddbegin \DIFadd{pAA}\DIFaddend . \cite{Loiseau2003a, LLauro2004} 
Given the difficulties of direct polymerization of AA, pAA is often made in a two-step process starting from a precursor monomer that is less affected 
by transfer reactions, followed by a polymer-analogous reaction that yields pAA. 

A canonical two-step method relies on the ester of acrylic acid and \emph{tert}-butanol, the latter acting analogously to a protecting group. 
\emph{tert}-Butyl acrylate (\DIFdelbegin %DIFDELCMD < {%%%
\DIFdel{t}%DIFDELCMD < }%%%
\DIFdelend \DIFaddbegin \emph{\DIFadd{t}}\DIFaddend BuAc) can be polymerized by anionic or radical means, and the resulting polymer is deesterified into pAA under acidic conditions,
trifluoroacetic acid in DCM (TFA/DCM) being especially common. \cite{Dulong2015, Yong2016, Chen2013a, Li2006, Li2008a, Lauber2017} 
However, even a strong molar excess yields pAA with at most 99.5\% liberated 
carboxylic acid, the rest being still esterified with 
\emph{tert}-butanol.\cite{Loiseau2003a, Chen2013a}  The 
incomplete cleavage remains unaddressed, even though residual \emph{tert}-butyl groups 
can affect polymer properties. Foreseeable effects are an increase of surface pressure at water-apolar interfaces and a greater tendency to 
form complexes through hydrophobic interactions, since the resulting structures are comparable to hydrophobically modified pAA. Thus, a means 
of obtaining ester-free pAA from p\emph{t}BuAc is desirable.

In the present work, we report a cleavage reaction that yields quantitatively de-\emph{
tert}-butylated pAA rapidly, based on the work of 
Palladino and co-workers\DIFdelbegin \DIFdel{\mbox{%DIFAUXCMD
\cite{Palladino2012}
}%DIFAUXCMD
.}\DIFdelend \DIFaddbegin \DIFadd{.\mbox{%DIFAUXCMD
\cite{Palladino2012}
}%DIFAUXCMD
}\DIFaddend We apply the method to the synthesis of a 
triblock copolymer with a monomer that is representative \cite{Lue2011, Schilli2004a}
of the use of pAA in literature, \DIFdelbegin \DIFdel{poly(}\emph{\DIFdel{N}}%DIFAUXCMD
\DIFdel{-iso}%DIFDELCMD < \-%%%
\DIFdel{propyl}%DIFDELCMD < \-%%%
\DIFdel{acryl}%DIFDELCMD < \-%%%
\DIFdel{amide) (pNIPAm)}\DIFdelend \DIFaddbegin \DIFadd{pNIPAm}\DIFaddend . Full 
\emph{tert}-butyl ester cleavage is shown
by NMR. To assert the necessity of complete deprotection of \DIFdelbegin \DIFdel{poly(}\DIFdelend \DIFaddbegin \DIFadd{p}\DIFaddend \emph{\DIFdelbegin \DIFdel{tert}\DIFdelend \DIFaddbegin \DIFadd{t}\DIFaddend }\DIFdelbegin \DIFdel{-butyl acrylate) }\DIFdelend \DIFaddbegin \DIFadd{BuAc }\DIFaddend precursors, we show that the ubiquitous 
TFA/DCM conditions yield polymers with markedly distinct flow and conformational properties, as witnessed by viscosity,
dynamic light scattering (DLS) and small-angle X-ray scattering measurements (SAXS).

\section*{Experimental section}
\textbf{Materials}
2,2'-Azobis(2-methylpropionitrile) (AIBN, 98\%), \emph{N}-iso\-propyl\-acryl\-amide (NIPAm,  97\%), \emph{tert}-butyl acrylate (\emph{t}BuAc, 
contains 10-20 ppm mono\-methyl ether hydro\-quinone as inhibitor, 98\%), trioxane ($\geqslant$99.9\%), tri\-fluoro\-acetic acid (TFA $\geqslant$99\%), 
concentrated HCl (37\% solution in water), and dioxane (anhydrous, 99.8\%) were obtained from Sigma-Aldrich, Germany, and used without further purification 
unless noted otherwise. NIPAm was recrystallized twice from a mixture of hexane and acetone. \emph{t}BuAc was passed over a short column of Al$_2$O$_3$ to
remove the inhibitor. AIBN was recrystallized thrice from MeOH. Bis(2-methylpropionic acid)trithiocarbonate \cite{Lai2002a, Cetintas2017} (BMAT) and
\DIFdelbegin \DIFdel{pNIPAm-BMAT }\DIFdelend \DIFaddbegin \DIFadd{BMAT-terminated pNIPAm }\DIFaddend \cite{Cetintas2017} were synthesized as described in literature\DIFaddbegin \DIFadd{, with NMR and SEC characterization described in Table \ref{Stab:orgsec}}\DIFaddend .
Hexafluoroisopropanol (HFIP, AR), dichloro\-methane (DCM, AR), tetra\-hydro\-furan (THF, HPLC-grade), and meth\-anol (MeOH, HPLC) were 
obtained from Biosolve, France.

\textbf{Synthesis of pNIPAm-\emph{b}-p\emph{t}\DIFdelbegin \DIFdel{BuAc-BMAT}\DIFdelend \DIFaddbegin \DIFadd{BuAc-}\emph{\DIFadd{b}}\DIFadd{-pNIPAm}\DIFaddend }
A Schlenk flask was charged with a solution of \DIFaddbegin \DIFadd{BMAT-terminated }\DIFaddend pNIPAm$_\mathrm{24}$ \DIFdelbegin \DIFdel{-BMAT }\DIFdelend (0.581 g, 0.104 mmol), \emph{t}BuAc (5.3 g, 41.4 mmol), 
AIBN (3.4 mg, 0.020 mmol) and trioxane (0.372 g, 4.14 mmol) in dioxane (21 mL), and was subjected
to five freeze-pump-thaw cycles. The solution was magnetically stirred at 70 $^\circ \mathrm{C}$ for 45 min, 
after which the polymerization was stopped by admitting atmosphere into the flask while cooling on an ice bath.
The material was precipitated thrice in ice-cold MeOH:water 3:1 and dried under high vacuum.

\textbf{TFA/DCM deprotection of pNIPAm-\emph{b}-p\emph{t}\DIFdelbegin \DIFdel{BuAc-BMAT}\DIFdelend \DIFaddbegin \DIFadd{BuAc-}\emph{\DIFadd{b}}\DIFadd{-pNIPAm}\DIFaddend }
pNIPAm-\emph{b}-p\emph{t}\DIFdelbegin \DIFdel{BuAc-BMAT }\DIFdelend \DIFaddbegin \DIFadd{BuAc-}\emph{\DIFadd{b}}\DIFadd{-pNIPAm }\DIFaddend (1.05 g, 0.023 mmol) was dissolved in DCM (50 mL) in a round-bottomed flask, 
to which was added 6 equivalents of TFA (4.2 g, 39 mmol) with respect to the amount of \emph{t}Bu units.
The mixture was stirred for 72 h at room temperature, the precipitated product was washed with fresh DCM,
and extensively dried under high vacuum. The dry polymer was dispersed in water by sonication and 
titrated with an NaOH solution until the solution was neutral. Water was removed by freeze drying.
This yielded \DIFdelbegin \DIFdel{the }\DIFdelend \DIFaddbegin \DIFadd{pNIPAm-}\DIFaddend (pAA-\emph{co}-p\emph{t}BuAc)-\emph{b}\DIFdelbegin \DIFdel{-pNIPAm-BMAT }\DIFdelend \DIFaddbegin \DIFadd{-pNIPAm }\DIFaddend as a white, fluffy powder.

\textbf{HCl/HFIP deprotection of pNIPAM-\emph{b}-p\emph{t}\DIFdelbegin \DIFdel{BuAc-BMAT}\DIFdelend \DIFaddbegin \DIFadd{BuAc-}\emph{\DIFadd{b}}\DIFadd{-pNIPAm}\DIFaddend }
pNIPAm-\emph{b}-p\emph{t}\DIFdelbegin \DIFdel{BuAc-BMAT }\DIFdelend \DIFaddbegin \DIFadd{BuAc-}\emph{\DIFadd{b}}\DIFadd{-pNIPAm }\DIFaddend (3.5 g, 0.093 mmol) was dissolved in HFIP 
(314 mL) in a round-bottomed flask, to which concentrated HCl 37\% (2.6 mL, 31.4 mmol) 
was added, 1.3 equivalents with respect to the amount of \emph{t}Bu units. After 4 h, 
the mixture was stripped of volatiles and the dry polymer was dispersed in 
water and titrated with an NaOH solution until the solution was neutral. A slight cloudiness was removed by
centrifugation at 15000g, and the supernatant was dehydrated by freeze drying. 
\DIFdelbegin \DIFdel{pAA-}\DIFdelend \DIFaddbegin \DIFadd{pNIPAm-}\DIFaddend \emph{b}\DIFdelbegin \DIFdel{-pNIPAm-BMAT }\DIFdelend \DIFaddbegin \DIFadd{-pAA-}\emph{\DIFadd{b}}\DIFadd{-pNIPAm }\DIFaddend was recovered as a white fluffy powder.

\textbf{Polymer characterization}
$^1$H- and $^{13}$C-NMR measurements were performed on a Bruker AVANCE 400 spectrometer. 
Fourier transform infrared absorption spectra were collected on a Bruker PLATINUM ATR mounted into a TENSOR
spectrometer on powdered sample. A background was taken and subtracted from the spectra, followed by
baseline correction and transformation from transmission to absorption.

\textbf{Size-exclusion chromatography}
Size exclusion chromatography (SEC) 
of pNIPAm and its copolymers with p\emph{t}BuAc was done on a Agilent 1200 
HPLC equipped with a PLgel $5 \mathrm{\mu m}$ 
Mixed-D column with RI detection, calibrated with \DIFdelbegin \DIFdel{poly(styrene) }\DIFdelend \DIFaddbegin \DIFadd{polystyrene }\DIFaddend standards, and HPLC-grade 
THF as eluent. The presence of pNIPAm blocks complicates the analyses due to 
column-analyte interactions. The problem is somewhat remedied \DIFdelbegin \DIFdel{by }\DIFdelend in an 
aqueous SEC system on corresponding de-\emph{tert}-butylated products. Aqueous size 
exclusion chromatography of pAA copolymers was done in an aqueous buffer of 
0.01 M Na$_2$HPO$_4$/NaH$_2$PO$_4$ with 0.1 M NaNO$_3$ on an Agilent 1260 Infinity 
II HPLC equipped with a Waters Ultrahydrogel500 column for molecular weight analysis.
\DIFaddbegin \DIFadd{Analytes were detected with a refractive index detector and by UV adsorption. 
}\DIFaddend 

\textbf{Trithiocarbonate cleavage with hydrazine}
In adaptation of a previous report,\cite{Shen2010}  
sodium salt of \DIFaddbegin \DIFadd{BMAT-derived }\DIFaddend triblock pNIPAm-\emph{b}\DIFdelbegin \DIFdel{-pAA-BMAT }\DIFdelend \DIFaddbegin \DIFadd{-pAA-}\emph{\DIFadd{b}}\DIFadd{-pNIPAm }\DIFaddend (0.100 g, 3.34 $\mathrm{\mu}$mol of
trithiocarbonates) was dissolved in water (11 mL), to which was added 5 equivalents of
hydrazine (1 $\mathrm{\mu}$L, 16.7 $\mathrm{\mu}$mol). The mixture was stirred at room temperature
for 72 h, and an aliquot was diluted 1:3 with 0.01 M Na$_2$HPO$_4$/NaH$_2$PO$_4$ with
0.1 M NaNO$_3$, which was analyzed with SEC over an Agilent PL AquaGel-OH 8 $\mathrm{\mu}$m column
and compared with an aliquot of non-treated sample.

\textbf{DLS of triblock solutions}
DLS measurements on triblock networks were performed in the manner described by
Bohdan et al.\cite{Bohdan2016} An ALV light scattering apparatus 
equipped with a 632.8 nm JDSU 1145P laser and a LSE-5004 correlator was used. Temperature was 
controlled by means of a Julabo water bath. DLS measures the electric field autocorrelation function 
$g_2(t) = {\langle I(0)I(t) \rangle}/{\langle I^2 \rangle}$, from which the normalized 
first-order correlation function $g_1(t)$ is obtained:
\begin{equation}
g_1(t) = \beta^{-1/2}\sqrt{g_2(t) - 1}
\end{equation}
in which $\beta$ is a geometry-dependent constant close to unity. We describe 
$g_1(t)$ by a two-mode stretched exponential decay:
\begin{equation} \label{eqn:bistrexp}
g_1(t) = \sum\limits_{i=f,s}A_i\exp(-\frac{t^\alpha_i}{\tau_i})
\end{equation}
where $\tau$ is the exponential decay time and $\alpha$ is the stretch parameter. The latter
characterises the width of the decay time distribution, with $\alpha_i=1$ corresponding to
a simple exponential decay and smaller values to a more stretched shape, and thus a broader
decay time distribution. The subscripts $f$ and $s$ refer to, respectively, the fast and
slow modes that we observe for all relaxations.\cite{Thuresson1999} 

Samples were made by addition of 1 mL water to 292 mg of polymer, and were filtrated over 
0.2 $\mathrm{\mu m}$ regenerated cellulose (RC) syringe filters after dissolution. 
Correlation curves were recorded for times between 3000 and 10000 s to ensure full 
decay of $g_1(t)$, with higher temperatures resulting in slower decays.

\textbf{Rheology of triblock solutions}
Rheological measurements were recorded on an Anton Paar MCR-301, using a 1 mL Couette geometry and temperature
control by a Peltier system. Evaporation was avoided by the use of a tetradecane-filled solvent trap. Viscosities
were measured by monitoring stress as a function of shear rate between 0.01 and 1000 $\mathrm{s}^{-1}$.

\textbf{Small-angle X-ray scattering of triblock solutions}
Small-angle X-ray scattering experiments were performed on a SAXSLAB Ganesha 300 XL set-up with a 
GeniX-Cu ultra low divergence source with a wavelength of 1.54 \AA$^{-1}$ and a flux of 
10$^8$ $s^{-1}$,
and a Pilatus 300K silicon pixel detector with 487x619 pixels of 172x172 $\mu \mathrm{m}^
2$. Temperature was controlled
between 20 and 60 $^\circ\mathrm{C}$ with a Julabo temperature controller. Three 
detector distances were used,
resulting in a $q$-range of 0.001 to 2 \AA$^{-1}$. Aqueous 292 mg mL$^{-1}$ 
samples were introduced into 
2 mm quartz capillaries, heated to 60 $^\circ\mathrm{C}$, equilibrated for 2 hours, then 
measured for 1800s at three different detector
distances. This was repeated for 40 and 20 $^\circ\mathrm{C}$. The diffraction patterns 
were azimuthally averaged
and intensities were corrected for transmission, sample thickness and detector distance. 
Finally, scattering from water-filled capillaries corrected in the same fashion was
subtracted from the scattering of triblock solutions.

\section*{Results and discussion}
We prepare pNIPAm-\emph{b}-pAA-\emph{b}-pNIPAm (NAN) triblock copolymers 
from pNIPAm-\emph{b}-\emph{t}BuAc-\emph{b}-pNIPAm (NTN) precursors, as 
shown in Figure \ref{syn-pNAN}, and further detailed in the Supporting Information 
(Figures \ref{Ssyn:pNBN}$^\dag$,\ref{Snmr:pNBN-short}$^\dag$). One common \emph{t}Bu-carrying precursor 
N$_\mathrm{24}$-T$_\mathrm{330}$-N$_\mathrm{24}$, where N$_\mathrm{24}$ denotes a 
pNIPAm block length of 24 monomers and $T_\mathrm{330}$ denotes a p\emph{t}BuAc 
block length of 330 monomers, serves to compare two acidic ester cleavage reactions: 
stirring with 6 equivalents of TFA in DCM for 72h and with 1.3 equivalents of HCl in HFIP for 4h. 
\DIFaddbegin \DIFadd{The times are sufficiently long to preclude any further reaction. For TFA/DCM, we choose
conditions similar to previous reports on pAA synthesis \mbox{%DIFAUXCMD
\cite{Dulong2015, Yong2016, Chen2013a, Li2006, Li2008a, Lauber2017}
}%DIFAUXCMD
to allow comparison to a situation that is typical in the field.
}\DIFaddend 

\begin{figure}
\centering
\includegraphics[width=0.95\linewidth]{syn-pNAN-complete.pdf}
\caption{Acidic de-\emph{tert}-butylation of the \DIFdelbeginFL \DIFdelFL{poly(}\emph{\DIFdelFL{tert}}%DIFAUXCMD
\DIFdelFL{-butyl acrylate) }\DIFdelendFL \DIFaddbeginFL \DIFaddFL{p}\emph{\DIFaddFL{t}}\DIFaddFL{BuAc }\DIFaddendFL block in an ABA 
copolymer with \DIFdelbeginFL \DIFdelFL{poly(}\emph{\DIFdelFL{N}}%DIFAUXCMD
\DIFdelFL{-isopropylacrylamide) }\DIFdelendFL \DIFaddbeginFL \DIFaddFL{pNIPAm }\DIFaddendFL using (left) TFA in DCM and (right) dilute 
stoichiometric HCl in HFIP. Characteristic degrees of \emph{t}-butoxy group removals are given as 
percentages. \label{syn-pNAN}} 
\end{figure}

NMR spectroscopy on aliquots from the reaction mixtures in MeOD shows
resonances characteristic of the NIPAm pendant groups at 1.16 and 3.97 ppm, as seen
in the NMR spectra of Figure \ref{nmr-pNAN}a. Contributions of the \emph{t}Bu peak appear
at 1.47 and 1.48 ppm, as the dominant peak in NTN precursor,
and as a well-noticeable contribution for the TFA-deprotected NAN, even after 72h. 
Contrarily, in the HCl-treated case, the residual \emph{t}Bu peak is sufficiently 
small to be indistinguishable from the tacticity features of the backbone protons, which 
are well-resolved in MeOD. \DIFaddbegin \DIFadd{Save for the resonances attributed to residual }\emph{\DIFadd{tert}}\DIFadd{-butyl,
the polymeric contributions to the spectra are identical, indicating no other changes
to the chemical structure than de-}\emph{\DIFadd{tert}}\DIFadd{-butylation.
}

\DIFaddend In Figure \ref{nmr-pNAN}b, the resonances of backbone
and \emph{t}Bu are enlarged, robustly confirming the disappearance of the latter in the 
case of HCl/HFIP after 3h with respect to the tacticity features. 
We note that the often-used D$_2$O does not provide an adequate environment for the 
\emph{t}Bu groups to give a sharp signal, and thus leads one to underestimate its presence. 

\begin{figure}
\centering
\includegraphics[width=\linewidth]{Spectral-data-final.pdf}
\caption{(a) $^\mathrm{1}$H-NMR spectra of MeOD solutions of a (top) \emph{t}BuAc 
copolymer with pNIPAm, 
(middle) product of deprotection with 6 eq. TFA in DCM, and (bottom) crude sample of deprotection 
mixture with 1.3 eq. HCl in HFIP. (b) Enlarged regions of crude samples in MeOD of deprotection with
TFA in DCM after 72h and with HCl in HFIP after 0.5h, 1h, 3h and 4h. Note that the HCl/HFIP spectra
are inflated eight times with respect to the TFA/DCM spectrum, after normalization to 
the pNIPAm methine
peak at 3.9 ppm. (c) Deconvolution of a spectrum taken on the deprotection in HCl/HFIP 
after 0.5h. The components corresponding to \emph{tert}-butyl protons are filled. \label{nmr-pNAN}}
\end{figure}

Deconvolution of the spectra between 1.4 and 1.8 ppm allows to separate the \emph{t}Bu 
contribution from the
backbone singlets by fitting each peak with a Gaussian-Lorentzian function. Figure \ref{nmr-pNAN}c 
shows a representative case. 
A comparison of the \emph{N}-isopropyl methine signal at 3.97 ppm to the \emph{t}Bu 
contribution at 1.47 ppm gives the deprotection efficiency $\phi_{D}$ according to

\begin{equation}
\phi_D=1-\frac{\frac{1}{9}A_{t\mathrm{Bu}}}{A_{Ni\mathrm{m}}}\frac{\mathrm{DP}_\mathrm{N}}{\mathrm{DP}_\mathrm{T}}
\end{equation}

where $A_{t\mathrm{Bu}}$ and ${A_{Ni\mathrm{m}}}$ are the integrals under the 
deconvoluted \emph{t}Bu contributions and the
\emph{N}-isopropyl methine signals, and $\mathrm{DP}$ is the degree of polymerization of 
the block denoted in subscript.
This procedure thus provides accurate values for $\phi_{D}$, even in cases where the area under the 
\emph{t}Bu region is dominated by backbone signals.

\DIFdelbegin \DIFdel{We speculate that the inferior ability of TFA/DCM to cleave }\emph{\DIFdel{t}}%DIFAUXCMD
\DIFdel{Bu esters and the 
rather high irreproducibility 
of $\phi_\mathrm{D, TFA/DCM}$ is related to inaccessibility of TFA to the reaction sites 
as the reaction nears completion. 
In the course of the first hour, the partially cleaved pAA-rich structures precipitate 
from the mixture and adhere to the reactor wall, likely preventing
an efficient access of TFA to the last few groups. In addition, carboxylates form dimers 
in apolar environments \mbox{%DIFAUXCMD
\cite{Fujii1988}
}%DIFAUXCMD
,
and therefore the TFA is possibly deactivated by complexation to the polymeric carboxyl 
groups that become progressively
available throughout the reaction. In contrast, HFIP dissolves all starting materials, 
intermediates and the product
of the deprotection. HCl is a non-carboxyl acid, and does not form carboxyl dimers. In 
short, both the improved solubility 
of the substrate and availability of de-}\emph{\DIFdel{tert}}%DIFAUXCMD
\DIFdel{-butylation agent favor HCl/HFIP as a
means to synthesize pAA from p}\emph{\DIFdel{t}}%DIFAUXCMD
\DIFdel{BuAc.
}%DIFDELCMD < 

%DIFDELCMD < %%%
\DIFdelend Given the improved rate and extent of the reaction, we strongly recommend deprotection 
with HCl/HFIP above TFA/DCM, with deprotection efficiencies of $>$99.9\% for HCl/HFIP 
and 90-98\% for TFA/DCM. We summarize the deprotection efficiencies 
$\phi_{D}$ of the two methods in Table \ref{tab:results}, $\phi_{D}$. The entries marked with 
an asterisk are used for a comparison of flow properties and DLS relaxation times 
(\emph{vide infra}). \DIFaddbegin \DIFadd{As seen in Table \ref{tab:results}, HCl/HFIP achieves
full }\emph{\DIFadd{tert}}\DIFadd{-butyl removal down to 1.0 equivalents of HCl. We note that the polymeric contributions
to the NMR spectra of all HCl/HFIP-derived products were identical for all acid loadings.
}\DIFaddend Additionally, Table \ref{tab:results} shows deprotection and aqueous SEC characterization data for runs on 
further NAN structures with slightly varied block lengths. FT-IR data are given 
in Supporting Information (Figure \ref{Sfig:ir}$^\dag$).

We use aqueous SEC to confirm the triblock identity of our compounds, and therefore the 
stability of the trithiocarbonate 
bond under conditions of deprotection and subsequent hydrolysis. We find similar values 
of $M_{n}$ for both TFA- and HCl-deprotected
species, as can be seen in Table \ref{tab:results}. Additionally, we intentionally 
cleaved a sample of N$_\mathrm{24}$-A$_\mathrm{330}$-N$_\mathrm{24}$
using hydrazine\DIFdelbegin \DIFdel{\mbox{%DIFAUXCMD
\cite{Shen2010}
}%DIFAUXCMD
.}\DIFdelend \DIFaddbegin \DIFadd{.\mbox{%DIFAUXCMD
\cite{Shen2010}
}%DIFAUXCMD
}\DIFaddend Treatment with hydrazine resulted in a discoloration of 
the sample and a reduction of $M_{n}$ by a factor
of two, as shown in Figure \ref{Sfig:aqs04}$^\dag$. Thus, the triblocks are quantitatively cleavable, 
which shows that the trithiocarbonate is stable, even after deprotection and neutralization to pH 7.

\DIFaddbegin \DIFadd{We argue that two factors improve the p}\emph{\DIFadd{t}}\DIFadd{BuAc de-esterification efficiency 
of HCl/HFIP with respect to TFA/DCM. DCM is not a solvent for pAA, which prevents
contact between the partially deprotected intermediates (essentially pAA/p}\emph{\DIFadd{t}}\DIFadd{BuAc 
copolymers) and the acidic medium. In addition, carboxylates form dimers 
in apolar environments,\mbox{%DIFAUXCMD
\cite{Fujii1988}
}%DIFAUXCMD
further decreasing contact.
Our results are not meant to capture the respective contributions of both factors, 
however, it is clear that TFA/DCM cannot be recommended for the synthesis of highly 
defined pAA. Because HFIP does not precipitate the product and since HCl is a non-carboxyl
acid, both issues are not at play for HCl/HFIP de-}\emph{\DIFadd{tert}}\DIFadd{-butylation.
}

\DIFaddend \begin{table}
\centering
\caption{Molecular weights and deprotection efficiencies of NAN copolymers. 
Block lengths are based on theoretical molecular weight $M_{th}$ that is based on 
monomer conversion. 
$M_n$ and \DIFdelbeginFL \DIFdelFL{$D$ }\DIFdelendFL \DIFaddbeginFL \DIFaddFL{$\DJ$ }\DIFaddendFL are calculated from SEC traces after elution in 
0.01 M aqueous phosphate buffer at pH 7 with 0.1 M NaNO$_3$. Entries marked with an 
asterisk are used further on in a comparative study of flow behavior, DLS relaxation times and SAXS.}
\begin{tabular}{lllllll}
\toprule
&Polymer structure 	&$M_{th}$ (kDa) & $M_{n}$  (kDa)&		\DIFdelbeginFL \DIFdelFL{$D$}\DIFdelendFL \DIFaddbeginFL \DIFaddFL{$\DJ$}\DIFaddendFL &		Treatment&	$\mathrm{\phi}_\mathrm{D}$\\
\midrule
&N$_\mathrm{26}$-(A$_\mathrm{312}$-T$_\mathrm{5}$)-N$_\mathrm{26}$	&	29.0	&	14.5		&	1.47		&TFA 6 eq	&	98.4\% \\
&N$_\mathrm{25}$-A$_\mathrm{362}$-N$_\mathrm{25}$							&	32.5	&	29.0	&	1.52	&HCl 1.2 eq&	$>$99.9\% \\
&N$_\mathrm{36}$-A$_\mathrm{220}$-N$_\mathrm{36}$			&	24.3	&	26.8		&	1.59	&HCl 1.0 eq 		&		$>$99.9\% \\
\textbf{*}&N$_\mathrm{24}$-(A$_\mathrm{297}$-T$_\mathrm{33}$)-N$_\mathrm{24}$	&	29.5	&	47.4		&	1.37		&TFA 6 eq	&	90.1\%	\\
\textbf{*}&N$_\mathrm{24}$-A$_\mathrm{330}$-N$_\mathrm{24}$							&	29.2	&	44.7		&	1.43		&HCl 1.3 eq&$>$99.9\%	\\
\bottomrule
\end{tabular}

\label{tab:results}
\end{table}

Literature abounds with reports on properties of pAA (co)polymers derived
from p\emph{t}BuAc\DIFdelbegin \DIFdel{\mbox{%DIFAUXCMD
\cite{Chen2013a, Maniego2017}
}%DIFAUXCMD
.}\DIFdelend \DIFaddbegin \DIFadd{.\mbox{%DIFAUXCMD
\cite{Chen2013a, Maniego2017}
}%DIFAUXCMD
}\DIFaddend We now turn to the question whether \emph{tert}-butyl moieties influence the solution 
properties of our copolymers strongly, even when present at the level we found for TFA/DCM. To this 
end, we studied the difference in the DLS relaxation time and viscosity of aqueous solutions above 
overlap between triblocks cleaved with TFA and HCl, derived from a common \DIFdelbegin \DIFdel{NBN }\DIFdelend \DIFaddbegin \DIFadd{NTN }\DIFaddend precursor, 
N$_\mathrm{26}$-T$_\mathrm{318}$-N$_\mathrm{26}$. A common NTN precursor ensures a constant 
chain length between the two polymers. The NTNs employed in the comparison are labelled with an asterisk in Table \ref{tab:results}.
In addition, SAXS was used to inquire into structural differences in HCl/HFIP and TFA/DCM products\DIFaddbegin \DIFadd{, respectively
N$_\mathrm{24}$-(A$_\mathrm{297}$-T$_\mathrm{33}$)-N$_\mathrm{24}$, where the center block in the latter is
the incompletely cleaved p}\emph{\DIFadd{t}}\DIFadd{BuAc center block}\DIFaddend . Additionally,
we compared two HCl-deprotected structures of different block lengths, N$_\mathrm{36}$-A$_\mathrm{220}$-N$_\mathrm{36}$
and N$_\mathrm{24}$-A$_\mathrm{330}$-N$_\mathrm{24}$. 

At a weight concentration of 292 mg mL$^{-1}$ without added salt, all polymers dissolved within two 
days on a shaking bench. Based on an estimate for $R_G$ of 5 nm taken from DLS measurements of 
dilute solutions at high salt (data not shown), the overlap concentration is below 
90 mg mL$^{-1}$ for \DIFdelbegin \DIFdel{both 
NANs }\DIFdelend \DIFaddbegin \DIFadd{all 
triblocks }\DIFaddend in demineralized water. Thus, our solutions strongly exceed overlap.
When heated to 60 $^\circ\mathrm{C}$, which is well above the bulk LCST for 
pNIPAm homopolymer around 32 $^\circ \mathrm{C}$, both solutions remain transparent. This is
in marked contrast to the rapid increase in turbidity that is observed for diblock solutions 
with a charged block and an LCST block, where the turbidity increase is attributed to 
formation of micelles due to the association
of pNIPAm blocks above the LCST. \cite{Schilli2004a, Chen2013a, Bayati2012} \DIFdelbegin \DIFdel{We }\DIFdelend \DIFaddbegin \DIFadd{In addition to the 
lack of a turibidty point, we }\DIFaddend note that our 
samples show no macroscopic phase separation, in contrast to solutions of pNIPAm homopolymer 
\DIFdelbegin \DIFdel{\mbox{%DIFAUXCMD
\cite{Heskins1968}
}%DIFAUXCMD
,copolymers\mbox{%DIFAUXCMD
\cite{Teodorescu2010a}
}%DIFAUXCMD
,}\DIFdelend \DIFaddbegin \DIFadd{,\mbox{%DIFAUXCMD
\cite{Heskins1968}
}%DIFAUXCMD
copolymers,\mbox{%DIFAUXCMD
\cite{Teodorescu2010a}
}%DIFAUXCMD
}\DIFaddend and derived hydrogels.\cite{Gan2010}
We remind the reader than pNIPAm and pAA form complexes at a pH far lower than of these samples. 
Indeed, the solubility of our structures is wholly consistent with the complete deprotonation 
of pAA at neutral pH.\cite{Schilli2004a} 

We note that the mode of association of \emph{t}Bu-free NAN triblocks above overlap
is \emph{not} related to micelle formation, as is reported for telechelic polymers 
with far longer stickers.\cite{DeGraaf2011, McCormick2008, Kirkland2008, Lin2013a} 
\DIFdelbegin \DIFdel{The }\DIFdelend \DIFaddbegin \DIFadd{Solutions of }\DIFaddend HCl-cleaved \DIFdelbegin \DIFdel{NAN solutions }\DIFdelend \DIFaddbegin \DIFadd{N$_\mathrm{24}$-A$_\mathrm{330}$-N$_\mathrm{24}$ }\DIFaddend remain transparent even after being 
heated to 60 $^\circ\mathrm{C}$ for three days, and show a decrease in scattering 
intensity at 90$^\circ$ when 
heated to 40 and 60 $^\circ \mathrm{C}$ (Figure \ref{Sdls:CRs}$^\dag$), precluding micellization.
Under identical conditions, the TFA-cleaved \DIFdelbegin \DIFdel{samples }\DIFdelend \DIFaddbegin \DIFadd{N$_\mathrm{24}$-(A$_\mathrm{297}$-T$_\mathrm{33}$)-N$_\mathrm{24}$ solutions
}\DIFaddend develop some turbidity and show occasional spikes in scattering intensity. \DIFaddbegin \DIFadd{However, SAXS measurements (}\emph{\DIFadd{vide infra}}\DIFadd{)
again preclude micellization as a cause. }\DIFaddend DLS measurements on dilute (5 mg mL$^{-1}$)
\DIFdelbegin \DIFdel{NAN }\DIFdelend \DIFaddbegin \DIFadd{triblock }\DIFaddend solutions show no changes in count rate or correlation times
upon heating from room temperature to 60 $^\circ\mathrm{C}$, and thus show no 
evidence of micellization, due to the very short sticker length and high osmotic pressure 
due to the charged blocks.

\begin{figure}
\centering
\begin{subfigure}[t]{.49\textwidth}
\includegraphics[width=.99\linewidth]{ABARHE02-LF37-LF38-viscosities-only-eta.pdf}
\end{subfigure}
\caption{\DIFdelbeginFL \DIFdelFL{(a) }\DIFdelendFL Shear rate-dependent viscosity of 292 mg mL$^{-1}$ solutions of
\DIFdelbeginFL \DIFdelFL{triblocks prepared
with TFA }\DIFdelendFL \DIFaddbeginFL \DIFaddFL{N$_\mathrm{24}$-}\DIFaddendFL (\DIFaddbeginFL \DIFaddFL{A$_\mathrm{297}$-T$_\mathrm{33}$)-N$_\mathrm{24}$ (}\DIFaddendFL squares) and 
\DIFdelbeginFL \DIFdelFL{HCl-derived }\DIFdelendFL \DIFaddbeginFL \DIFaddFL{N$_\mathrm{24}$-A$_\mathrm{330}$-N$_\mathrm{24}$ }\DIFaddendFL (circles) at 15 $^\circ\mathrm{C}$ 
(blue) and 60 $^\circ\mathrm{C}$ (red).
\label{fig:rheo}}
\end{figure}

When heated from 15 to 60 $^\circ\mathrm{C}$, a 292 mg mL$^{-1}$ solution of 
\DIFdelbegin \DIFdel{the HCl batch }\DIFdelend \DIFaddbegin \DIFadd{N$_\mathrm{24}$-A$_\mathrm{330}$-N$_\mathrm{24}$ }\DIFaddend shows a strong increase in viscosity, 
whereas none is visually distinguishable \DIFdelbegin \DIFdel{in the TFA batch}\DIFdelend \DIFaddbegin \DIFadd{for solutions of 
N$_\mathrm{24}$-(A$_\mathrm{297}$-T$_\mathrm{33}$)-N$_\mathrm{24}$}\DIFaddend . 
This behavior is quantified in Figure \ref{fig:rheo} 
(left) by means of rotational rheology in a Couette geometry, where we plot the shear-rate dependent 
viscosity for both \DIFdelbegin \DIFdel{HCl and TFA-derived NAN at identical degree of polymerization, at temperatures significantly lower and significantly higher than the LCST}\DIFdelend \DIFaddbegin \DIFadd{triblocks, at $T$ above and below the bulk LCST of pNIPAm homopolymer}\DIFaddend .
All viscosity curves are characterized by a wide Newtonian range, and the Newtonian 
viscosity of the sample increased 70-fold \DIFdelbegin \DIFdel{in the HCl batch, 
whereas for the case of TFA the }\DIFdelend \DIFaddbegin \DIFadd{for N$_\mathrm{24}$-A$_\mathrm{330}$-N$_\mathrm{24}$, 
whereas the }\DIFaddend increase is marginal \DIFaddbegin \DIFadd{for the case  of
N$_\mathrm{24}$-(A$_\mathrm{297}$-T$_\mathrm{33}$)-N$_\mathrm{24}$}\DIFaddend . 

We stress that our solutions, upon heating, neither undergo a macroscopic phase 
transition nor have a cloud point. We merely perceive the effect of the collapse of pNIPAm
\DIFdelbegin \DIFdel{above 
its LCST through observing }\DIFdelend \DIFaddbegin \DIFadd{as }\DIFaddend the slowed-down relaxation of a transient network. Exactly 
the unchanged transparency and volume of the samples allowed us to 
employ DLS to study differences in 
relaxation times between the two solutions above overlap, below and above the bulk pNIPAm LCST. 
In Figure \ref{fig:dls}(a), we show averaged first-order normalized correlation 
functions $g_1(t)$
for both \DIFdelbegin \DIFdel{batches }\DIFdelend \DIFaddbegin \DIFadd{triblocks }\DIFaddend at 20, 40 and 60 $^\circ\mathrm{C}$. At all temperatures, 
the correlation functions eventually decay to zero after sufficiently 
long measurement times (up to 10000 s).

\begin{figure}
\begin{subfigure}[t]{.49\textwidth}
\includegraphics[width=.99\linewidth]{ABADLS03-normalized-correlations-fit.pdf}
\end{subfigure}%
\begin{subfigure}[t]{.49\textwidth}
\includegraphics[width=.99\linewidth]{Tau-vs-T-from-stretched-exp.pdf}
\end{subfigure}
\caption{(a) First-order correlation functions from DLS on 292 mg mL$^{-1}$ solutions of 
\DIFdelbeginFL \DIFdelFL{the same triblocks
}\DIFdelendFL \DIFaddbeginFL \DIFaddFL{N$_\mathrm{24}$-A$_\mathrm{330}$-N$_\mathrm{24}$ (squares) and N$_\mathrm{24}$-(A$_\mathrm{297}$-T$_\mathrm{33}$)-N$_\mathrm{24}$
(circles) }\DIFaddendFL at 20 $^\circ\mathrm{C}$ (blue), 40 $^\circ\mathrm{C}$ (orange), and 60 $^\circ\mathrm{C}$ (red). The lines
are fitted using Equation \ref{eqn:bistrexp}. (b) Characteristic slow and fast relaxation times $\tau_s$ and $\tau_f$ as 
function of temperature as fitted using Equation \ref{eqn:bistrexp}. The lines are \DIFdelbeginFL \DIFdelFL{meant }\DIFdelendFL \DIFaddbeginFL \DIFaddFL{intended }\DIFaddendFL as a guide to the eye.
\label{fig:dls}}
\end{figure}

All measured correlation functions are well described by a sum of two stretched exponentials
(Equation \ref{eqn:bistrexp}), and thus their relaxation is characterized by one slow and one 
fast mean relaxation time $\tau_\mathrm{s}$ and $\tau_\mathrm{f}$. The fits are plotted as 
solid lines in Figure \ref{fig:dls}a, and the relaxation times are plotted against temperature in 
Figure \ref{fig:dls}b. A third mode in $g_1(t)$ is also seen, but the acquisition of data of sufficient
statistics is prohibitively long for samples at high $T$, as is clear from 
Figure \ref{fig:dls}a. Since its contribution is only minor, it is neglected in the analysis.

The stretch exponents of the fast mode $\alpha_f$ are very close to unity, and the fast 
modes are thus exponential\DIFdelbegin \DIFdel{, and }\DIFdelend \DIFaddbegin \DIFadd{. Thus, }\DIFaddend the fast relaxation times are narrowly distributed. Most 
of the fast relaxations contribute very little to the overall function, and
only \DIFdelbegin \DIFdel{for the }\DIFdelend \DIFaddbegin \DIFadd{in solutions of N$_\mathrm{24}$-(A$_\mathrm{297}$-T$_\mathrm{33}$)-N$_\mathrm{24}$ 
at }\DIFaddend 20 $^\circ\mathrm{C}$ \DIFdelbegin \DIFdel{TFA-cleaved solution }\DIFdelend does the fast mode contribute more than 1\% of the signal.

The slow modes exhibit varying degrees of stretching with $\alpha_s$ between 0.76 and 0.95, 
$\alpha_s$ increasing with temperature. This suggests that the relaxation rates are more narrowly
distributed for transient networks in which the associators have a higher tendency to stick.
The dependence of $\tau_s$ mirrors the chemistry-dependent thermothickening observed in viscometry.
Both the absolute value of $\tau_s$ and their temperature variation mode are strongly affected 
by the chemistry of the middle block. For \DIFdelbegin \DIFdel{the HCl-cleaved NAN}\DIFdelend \DIFaddbegin \DIFadd{N$_\mathrm{24}$-A$_\mathrm{330}$-N$_\mathrm{24}$}\DIFaddend , 
$\tau_s$ is always longer, and increases over more than two orders of magnitude. Its 
TFA-cleaved counterpart has a more rapid relaxation of the slow mode, and shows only 
modest growth with temperature. 

\begin{figure}
\centering
\includegraphics[width=.80\linewidth]{Transient_network_d_tau.pdf}
\caption{Schematic of two snapshots of a cluster of a few chains in a transient network formed 
by ABA triblock copolymers. Note that in the range studied, temperature and pAA purity do not change 
the qualitative features of the system: they change only $\tau$, which is the waiting time between 
the snapshots, and is chosen to be above the node residence time and the decorrelation time of the 
structure on lengthscales on the order of this box size. The configuration at $t_0$ is drawn in the 
rightmost figure as dashed lines. $d_\mathrm{node}$ is the characteristic internode distance 
(\emph{vide infra} for its measurement by SAXS). For clarity, the nodes have been portrayed as 
stationary. \label{fig:network}}
\end{figure}

Three-mode stretched exponential relaxation of $g_1(t)$ in DLS measurements of transient polymer
networks has been previously described by other authors\DIFdelbegin \DIFdel{\mbox{%DIFAUXCMD
\cite{Bohdan2016, Thuresson1999}
}%DIFAUXCMD
, }\DIFdelend \DIFaddbegin \DIFadd{, \mbox{%DIFAUXCMD
\cite{Bohdan2016, Thuresson1999}
}%DIFAUXCMD
}\DIFaddend where poly(ethylene glycol) networks were formed by association of aliphatic chains, which 
were attached to the polymers at either one or both ends, yielding ``diblock'' or ``triblock'' 
structures. The relaxation 
times of the second mode increased strongly with increasing proportion of 
triblock structures (thus increasing network connectivity), and is ascribed to increased 
node residence times. Thus, this connectivity increase is analogous to a temperature 
increase in our experiment.
The $\tau$ of the third mode is reported to show an even stronger connectivity dependence
\cite{Bohdan2016} and is ascribed to diffusion of associated clusters. 


Similarly, we attribute the increase of $\eta$ and $\tau_s$ with \DIFdelbegin \DIFdel{temperature
}\DIFdelend \DIFaddbegin \DIFadd{$T$ 
}\DIFaddend to a prolongation of the mean residence time of pNIPAm stickers in the nodes, 
driven by a higher segregation tendency. The collective effect of the increased 
residence time is a slowing down of the transient network. The wavevector $q$ is 
$0.019\mathrm{nm}^{-1}$, which means that we 
are at length scales one order of magnitude beyond the node-node correlations.
Nonetheless, while $\tau_s$ cannot be trivially converted to a residence time,
the relaxation time must necessarily increase with residence time.\cite{Bohdan2016} 
We convey the relaxation of these transient networks in Figure 
\ref{fig:network} --- one must realize that the temperature and the extent of 
deprotection determine the $\tau$ over which we would observe relaxation.


It is remarkable that the availability of hydrophobic groups \emph{decreases} the residence time, 
since it could be argued that the hydrophobic effect would be stronger with the participation of such groups, 
and thus sticker association more prominent. However, we speculate that the residual esters 
actually improve compatibility between the pNIPAm and the pAA blocks, 
resulting in \emph{weakening} the physical crosslinks. Additionally, it could 
be argued that the coils slightly shrink causing them to overlap less. This is 
consistent with the lower viscosity of the \DIFdelbegin \DIFdel{TFA-derived }\DIFdelend \DIFaddbegin \DIFadd{N$_\mathrm{24}$-(A$_\mathrm{297}$-T$_\mathrm{33}$)-N$_\mathrm{24}$
}\DIFaddend solution at low temperatures, and more prominent participation of a fast mode at room temperature.


Additional clues for the origin of short sticker residence times for identical pNIPAm moieties
connected to more strongly hydrophobically contaminated pAA come from SAXS, which 
reveals time-averaged 
density fluctuations, thus structural information rather than $\tau$-dependent observables.
Here, we compare the angle-dependent absolute scattering intensities $\frac{d\Sigma}{d\Omega}(q)$ 
of 292 mg mL$^{-1}$ solutions of 
the \emph{tert}-butylated N$_\mathrm{24}$-(A$_\mathrm{297}$-T$_\mathrm{33}$)-N$_\mathrm{24}$ 
and fully liberated triblocks 
N$_\mathrm{36}$-A$_\mathrm{220}$-N$_\mathrm{36}$ and 
N$_\mathrm{24}$-A$_\mathrm{330}$-N$_\mathrm{24}$ at
20, 40 and 60 $^\circ\mathrm{C}$. All scattering intensities are peaked in the 
vicinity of $q^*=$0.03 \AA$^{-1}$, as is
visible in Figure \ref{fig:saxs}. The peaks sharpen and shift to smaller $q$ with $T$,
which suggests their origin to be density fluctuations due to increased non-solubility 
of pNIPAm with a corresponding internode distance $d_\mathrm{node}$.
Indeed, a change occurs between 20 and 40 $^\circ\mathrm{C}$, in agreement with an 
LCST of 32$^\circ\mathrm{C}$. We attribute the presence of an internode peak below the 
LCST to pNIPAm-pNIPAm interactions
that have been observed in scattering studies of aqueous pNIPAm before\DIFdelbegin \DIFdel{\mbox{%DIFAUXCMD
\cite{Lang2017}
}%DIFAUXCMD
,}\DIFdelend \DIFaddbegin \DIFadd{,\mbox{%DIFAUXCMD
\cite{Lang2017}
}%DIFAUXCMD
}\DIFaddend which note a decrease in excluded 
volume starting already from 5$^\circ\mathrm{C}$, and a consequent clustering below the 
LCST.\cite{Kawaguchi2009} 

Upon heating from 20 to 60 $^\circ\mathrm{C}$, the internode distance $d_\mathrm{node}$ that 
correspond to $q*$, at which $\frac{d\Sigma}{d\Omega}$ passes through a prominent maximum, 
increases from 18.5 nm to 20.4 nm for N$_\mathrm{36}$-A$_\mathrm{220}$-N$_\mathrm{36}$. Thus,
the pAA bridges are stretched to approximately twice their diameter of 10 nm, costing
approximately one $k_BT$ assuming entropic elasticity. For  
N$_\mathrm{24}$-A$_\mathrm{330}$-N$_\mathrm{24}$ $d_\mathrm{node}$ is between 19.3 nm and 21.9 nm,
in line with its longer pAA block. In other regards, the six curves for HCl-derived NAN 
structures are highly similar, suggesting that the qualitative picture in Figure \ref{fig:network}
is correct for all cases, and that the chemistry and temperature merely \DIFdelbegin \DIFdel{causes }\DIFdelend \DIFaddbegin \DIFadd{cause }\DIFaddend changes
in the internode distance and sticker residence time of the transient network.

\begin{figure}
\centering
\includegraphics[width=.65\linewidth]{abasax_prcsd_lf37-lf38-ivh-cleaned.pdf}
\caption{
X-ray scattering intensity as function of modulus of scattering
vector $q$ for 292 mg mL$^{-1}$ solutions of one TFA-deprotected and two 
HCl-deprotected NAN triblock copolymers 
at 20, 40 and 60 $^\circ\mathrm{C}$ in respectively blue, orange and red (from bottom to top, 
since the scattering decreases slightly with $T$). The labels $\mathrm{x=}n$ give the scaling factor
for the curves, used for better readability. The lines are plotted from the full data, while the
marks are plotted only every sixth point.
\label{fig:saxs}}
\end{figure}

The TFA-deprotected N$_\mathrm{24}$-(A$_\mathrm{297}$-T$_\mathrm{33}$)-N$_\mathrm{24}$ 
shows smaller values of $d_\mathrm{node}$, between 17.6 nm and 19.7 nm, corresponding to a lesser
stretching of the polyelectrolyte block due to more compatibility with the pNIPAm 
moieties. A striking feature of the curves is the high-$q$ region, where all curves show power-
laws of, in the case of
\DIFdelbegin \DIFdel{TFA-derived samples}\DIFdelend \DIFaddbegin \DIFadd{N$_\mathrm{24}$-(A$_\mathrm{297}$-T$_\mathrm{33}$)-N$_\mathrm{24}$}\DIFaddend , $q^{-2/3}$, and, 
in the case of \DIFdelbegin \DIFdel{HCl-derived }\DIFdelend \DIFaddbegin \DIFadd{fully liberated }\DIFaddend samples, $q^{-5/3}$, the latter 
corresponding to a self-avoiding walk\DIFdelbegin \DIFdel{\mbox{%DIFAUXCMD
\cite{DeGennes1979}
}%DIFAUXCMD
.}\DIFdelend \DIFaddbegin \DIFadd{.\mbox{%DIFAUXCMD
\cite{DeGennes1979}
}%DIFAUXCMD
}\DIFaddend We observe these forms up to a 
$d$ of 1.5 nm, and starting from 5.4 nm for TFA and around 2 nm 
for the HCl-derived samples, where scattering is dominated by chain statistics ($R_{g}>d>l_k$). 

At present, we have no interpretation of the $q^{-2/3}$ scaling for the TFA-treated case. Clearly,
it is not reasonable to attribute it to an even further stretching of the chains than we would
expect for a rod ($q^{-1}$). Thus, the observed shape is likely due to a superposition of different
dependencies in the high-$q$ scattering. Nonetheless, 
we stress that the high-$q$ data witness that TFA leaves us with 
pAA blocks that have strongly altered conformational
properties with respect to pure pAA blocks.
Whereas the structure on the internode length scale is only weakly dependent on 
temperature and chemistry,
the density fluctuations on the chain length scale are strongly altered by \emph{tert}-butyl 
moieties incompletely cleaved by TFA.

% or here conclusion?
Through DLS, rheology and visual observations of NAN triblocks around RT and at
60 $^\circ\mathrm{C}$, we observe a significant enhancement of sticker residence times 
of pNIPAm blocks 
when the \emph{tert}-butyl group is removed completely, whereas association seems to be minimal
when hydrophobic groups have presence. SAXS reveals internode distances of around twice the
dilute coil diameter that increase slightly with pAA chain length and absence of 
hydrophobic groups. More strikingly,
the density fluctuations on the chain length scales are dramatically different in the 
case when \emph{tert}-butyl groups are
not completely removed. Thus, the presence of these groups has a strong effect on the 
self-assembly behavior of NAN copolymers, 
which are prepared throughout literature through a method that we show does \emph{not} 
remove \emph{tert}-butyl up to an 
extent that their influence is negligible.

\section*{Conclusion}
In summary, we report a facile, rapid de-\emph{tert}-butylation scheme for the synthesis 
of pAA block copolymers 
with pNIPAm from the corresponding copolymer with p\emph{t}BuAc. The action of 
stoichiometric HCl on a HFIP solution
of the \emph{tert}-butylated precursor yields pAA blocks of a chemical purity that has been 
inaccessible \emph{via} the commonly used TFA/DCM mixture. In addition, we show that the 
presence of residual
\emph{tert}-butyl moieties, at a degree tolerated by the latter method, impacts the 
aqueous self-assembly properties 
of these polymers to such an extent that the thermothickening effect due to LCST-driven 
pNIPAm association is
essentially disabled for polymer solutions significantly above overlap. Thus, HCl/HFIP 
is an essential tool in the
study of the aqueous association of pAA copolymers derived from their \emph{tert}-butylated 
precursors.

\section*{Acknowledgements}
This work was financially supported by the Dutch Science Foundation (NWO VIDI Grant 723.013.005).
IKV thanks the Dutch Science Foundation (NWO ECHO Grant No. 712.016.002, NWO VIDI Grant 723.014.006) for
financial support. The authors wish to thank Jasper van der Gucht, Thomas Kodger and Dmitri Filippov
for their useful comments.

\DIFaddbegin \section*{\DIFadd{Supporting Information}}
\DIFadd{Table S1, Figures S1-S4: Detailed chemical polymer characterization. Figure S5: Light scattering intensities
of polymer solutions.
}

\DIFaddend \bibliography{library} %You need to replace "rsc" on this line with the name of your .bib file
\bibliographystyle{ieeetr} %the RSC's .bst file
\DIFaddbegin 


\newpage
\thispagestyle{empty}
\includegraphics[width=8.3cm,height=3.5cm]{Tocart.png} \\
\DIFadd{For Table of Contents Use Only }\\
\DIFadd{Title: ''Rapid and quantitative de-tert-butylation for poly(acrylic acid) block copolymers and influence on relaxation of thermoassociated transient network`` }\\
\DIFadd{Author(s): Filippov, Alexei; van Hees, Ilse; Fokkink, Remco; Voets, Ilja; Kamperman, Marleen  }\\
\DIFaddend 

\end{document}
